\section{Part I (b): Coupling Through Axes}

\subsection{Coupling}

The framework's central discovery is that scalings and rotations, though independent
by construction in the scalar case, are coupled when the scale is directional. A scaling
is a pure stretch—it changes lengths and rotates nothing—so as a matrix operator it is
symmetric. A rotation generator changes no length, so it is antisymmetric. The coupling
between them follows from these definitions alone, forced by the algebra of symmetric
and antisymmetric matrices.

\begin{definition}
A \textbf{scaling} is a symmetric matrix.
\end{definition}

\begin{definition}
A \textbf{rotation generator} is an antisymmetric matrix.
\end{definition}

\begin{definition}
The \textbf{bracket} of two operators $S$ and $T$ is their commutator:
\[\text{br}(S, T) = ST - TS.\]
\lean{Coupling.br}
\end{definition}

\begin{theorem}[The Coupling]
Two scalings bracket to a rotation. If $S$ and $T$ are both symmetric matrices, then
their bracket $\text{br}(S, T)$ is antisymmetric.
\end{theorem}
\begin{proof}
Let $S$ and $T$ be symmetric, so $S^T = S$ and $T^T = T$. Then
\[\text{br}(S, T)^T = (ST - TS)^T = T^T S^T - S^T T^T = TS - ST = -\text{br}(S, T),\]
which is the definition of antisymmetry. The coupling follows from transposition alone,
independent of any metric or coordinate choice.
\end{proof}
\lean{Coupling.bracket\_scaling\_scaling}

\begin{theorem}
A rotation and a scaling bracket to a scaling. If $A$ is antisymmetric and $S$ is symmetric,
then $\text{br}(A, S)$ is symmetric.
\end{theorem}
\begin{proof}
Let $A$ be antisymmetric and $S$ symmetric. Then
\[\text{br}(A, S)^T = (AS - SA)^T = S^T A^T - A^T S^T = SA - AS = -(AS - SA) + 2SA = \text{br}(A, S) - 2(AS - SA) + 2SA.\]
More directly:
\[\text{br}(A, S)^T = (AS - SA)^T = S^T A^T - A^T S^T = (-A)S - S(-A) = -AS + SA = -(AS - SA) = \text{br}(A,S),\]
using $A^T = -A$ and $S^T = S$.
\end{proof}
\lean{Coupling.bracket\_rotation\_scaling}

\begin{theorem}
Rotations close under the bracket. If $A$ and $B$ are both antisymmetric, then $\text{br}(A, B)$ is antisymmetric.
\end{theorem}
\begin{proof}
Let $A$ and $B$ be antisymmetric. Then
\[\text{br}(A, B)^T = (AB - BA)^T = B^T A^T - A^T B^T = (-B)(-A) - (-A)(-B) = BA - AB = -\text{br}(A, B),\]
which is antisymmetric.
\end{proof}
\lean{Coupling.bracket\_rotation\_rotation}

\begin{theorem}[The Coupling Structure]
The scale sector and rotation sector form a Cartan decomposition: $[S, S] \subseteq A$,
$[A, S] \subseteq S$, and $[A, A] \subseteq A$. Explicitly, if $S, T$ are symmetric and
$A, B$ are antisymmetric, then $\text{br}(S, T)$ is antisymmetric, $\text{br}(A, S)$ is symmetric, and $\text{br}(A, B)$ is antisymmetric.
\end{theorem}
\begin{proof}
This is the immediate consequence of the three preceding theorems.
\end{proof}
\lean{Coupling.coupling\_structure}

\begin{theorem}[Single Direction Generates Nothing]
The bracket of an operator with itself vanishes: $\text{br}(S, S) = 0$ for any matrix $S$.
\end{theorem}
\begin{proof}
$\text{br}(S, S) = SS - SS = 0$ by definition.
\end{proof}
\lean{Coupling.bracket\_self}

\begin{theorem}
Commuting matrices bracket to zero. If $S$ and $T$ satisfy $ST = TS$, then $\text{br}(S, T) = 0$.
\end{theorem}
\begin{proof}
$\text{br}(S, T) = ST - TS = 0$ when $ST = TS$.
\end{proof}
\lean{Coupling.no\_rotation\_of\_commuting}

\begin{definition}
In three spatial dimensions, the scaling in the $0$--$1$ plane is the matrix
\[K_1 = \begin{pmatrix} 0 & 1 & 0 \\ 1 & 0 & 0 \\ 0 & 0 & 0 \end{pmatrix}.\]
\end{definition}
\lean{Coupling.K$_1$}

\begin{definition}
The scaling in the $0$--$2$ plane is
\[K_2 = \begin{pmatrix} 0 & 0 & 1 \\ 0 & 0 & 0 \\ 1 & 0 & 0 \end{pmatrix}.\]
\end{definition}
\lean{Coupling.K$_2$}

\begin{definition}
The rotation in the $1$--$2$ plane is
\[J = \begin{pmatrix} 0 & 0 & 0 \\ 0 & 0 & 1 \\ 0 & -1 & 0 \end{pmatrix}.\]
\end{definition}
\lean{Coupling.J}

\begin{theorem}[Thomas–Wigner Rotation]
The bracket of the two scalings equals the rotation: $\text{br}(K_1, K_2) = J$.
\end{theorem}
\begin{proof}
Direct computation of the matrix product and subtraction.
\end{proof}
\lean{Coupling.boost\_bracket\_eq\_rotation}

\begin{theorem}
The two scalings do not commute: $K_1 K_2 \neq K_2 K_1$.
\end{theorem}
\begin{proof}
Suppose they commute. Then $\text{br}(K_1, K_2) = 0$. But $\text{br}(K_1, K_2) = J$, and the $(1, 2)$ entry of $J$ is $1 \neq 0$, a contradiction.
\end{proof}
\lean{Coupling.scalings\_do\_not\_commute}

\subsection{Algebra}

The coupling forces the framework to determine which Lie algebra it is. A maximal abelian
subalgebra of the scaling sector has dimension equal to the real rank of the Lorentz algebra.
Two independent results—the drift is a single linear functional, and commuting scalings
generate no rotation—imply there is exactly one scale direction specifiable simultaneously.
This forces the algebra to have real rank one, selecting it uniquely from the classified list.

\begin{theorem}[Orthogonality Under the Trace Form]
Scalings and rotations are orthogonal under the trace form on matrices. If $S$ is symmetric
and $A$ is antisymmetric, then $\text{tr}(SA) = 0$.
\end{theorem}
\begin{proof}
We have
\[\text{tr}((SA)^T) = \text{tr}(S),\]
the standard property of the trace. But $(SA)^T = A^T S^T = (-A) S = -AS$, so
\[\text{tr}(SA) = \text{tr}((SA)^T) = \text{tr}(-AS) = -\text{tr}(AS) = -\text{tr}(SA),\]
using the cyclic property $\text{tr}(XY) = \text{tr}(YX)$. Therefore $2 \text{tr}(SA) = 0$, giving $\text{tr}(SA) = 0$.
\end{proof}
\lean{Algebra.scaling\_rotation\_orthogonal}

\begin{definition}
A boost along direction $v \in \mathbb{R}^k$ is the $(k+1) \times (k+1)$ matrix
\[\text{boost}(v) = \begin{pmatrix} 0 & v^T \\ v & 0 \end{pmatrix},\]
where the top-left entry is zero, the top-right is the row vector $v^T$, the bottom-left
is the column vector $v$, and the bottom-right block is the zero $k \times k$ matrix.
\end{definition}
\lean{Algebra.boost}

\begin{theorem}
The $(0, 0)$ entry of a boost is zero.
\end{theorem}
\begin{proof}
Immediate from the definition.
\end{proof}
\lean{Algebra.boost\_zero\_zero}

\begin{theorem}
For a boost and spatial index $j$, the $(0, j+1)$ entry is $v_j$.
\end{theorem}
\begin{proof}
Immediate from the definition.
\end{proof}
\lean{Algebra.boost\_zero\_succ}

\begin{theorem}
For a boost and spatial index $i$, the $(i+1, 0)$ entry is $v_i$.
\end{theorem}
\begin{proof}
Immediate from the definition.
\end{proof}
\lean{Algebra.boost\_succ\_zero}

\begin{theorem}
For a boost and spatial indices $i, j$, the $(i+1, j+1)$ entry is zero.
\end{theorem}
\begin{proof}
Immediate from the definition.
\end{proof}
\lean{Algebra.boost\_succ\_succ}

\begin{theorem}
A boost is a scaling: the matrix $\text{boost}(v)$ is symmetric.
\end{theorem}
\begin{proof}
The transpose swaps the top-right and bottom-left blocks, which are $v^T$ and $v$ respectively.
This swaps a row vector with a column vector, but when tranposed they are identical, so
$\text{boost}(v)^T = \text{boost}(v)$.
\end{proof}
\lean{Algebra.boost\_isScaling}

\begin{theorem}
Boosts superpose: $\text{boost}(v + w) = \text{boost}(v) + \text{boost}(w)$.
\end{theorem}
\begin{proof}
This follows immediately from the linearity of the boost construction in its argument.
\end{proof}
\lean{Algebra.boost\_add}

\begin{theorem}
Boosts scale: $\text{boost}(c \cdot v) = c \cdot \text{boost}(v)$ for any scalar $c$.
\end{theorem}
\begin{proof}
Again, linearity of the construction.
\end{proof}
\lean{Algebra.boost\_smul}

\begin{theorem}
For $v, w \in \mathbb{R}^k$ and spatial indices $i,j$, the spatial block of the
product of two boosts is
\[
(\text{boost}(v)\,\text{boost}(w))_{i+1,j+1} = v_i w_j.
\]
\end{theorem}
\begin{proof}
Directly from matrix multiplication: the $(i{+}1,j{+}1)$ entry sums over the
shared index, and only the $0$-th (time) component of each boost is
nonzero off the diagonal block, so the sum collapses to the single term
$v_i w_j$.
\end{proof}
\lean{Algebra.boost\_mul\_succ\_succ}

\begin{theorem}[The Bracket of Two Boosts]
The bracket of two boosts, restricted to the spatial block, is the wedge of their
direction vectors. For $v, w \in \mathbb{R}^k$ and $i, j \in \{0, \ldots, k-1\}$,
\[[\text{boost}(v), \text{boost}(w)]_{i+1, j+1} = v_i w_j - v_j w_i.\]
This is the Thomas--Wigner rotation made explicit.
\end{theorem}
\begin{proof}
The computation is direct: the spatial block of $\text{boost}(v) \text{boost}(w)$ is $vv^T w^T + \ldots$,
and the antisymmetric part under subtraction yields $v_i w_j - v_j w_i$.
\end{proof}
\lean{Algebra.boost\_bracket\_spatial}

\begin{theorem}[Rank One Condition]
Two boosts commute if and only if their direction vectors are componentwise parallel:
$\text{boost}(v) \text{boost}(w) = \text{boost}(w) \text{boost}(v)$ if and only if
$v_i w_j = v_j w_i$ for all $i, j$.
\end{theorem}
\begin{proof}
If the boosts commute, their bracket vanishes. By the previous theorem, this means $v_i w_j - v_j w_i = 0$
for all $i, j$. Conversely, if $v_i w_j = v_j w_i$ for all $i, j$, then the bracket vanishes by the same calculation,
so the boosts commute.
\end{proof}
\lean{Algebra.boosts\_commute\_iff\_parallel}

\begin{theorem}[Maximal Abelian Subalgebra is One-Dimensional]
All simultaneously specifiable scale directions lie on one line. If $u$ has a nonzero component
$u_{i_0} \neq 0$ and $\text{boost}(u)$ commutes with $\text{boost}(v)$, then
$v = \frac{v_{i_0}}{u_{i_0}} u.$
\end{theorem}
\begin{proof}
Since the boosts commute, we have $u_i w_j = u_j w_i$ for all $i, j$ by the previous theorem.
Rearranging with $i = i_0$ gives $w_i = \frac{w_{i_0}}{u_{i_0}} u_i$ for all $i$.
\end{proof}
\lean{Algebra.commuting\_boosts\_lie\_on\_a\_line}

\begin{theorem}
A commuting family of scalings is a one-parameter family. If $u$ is a direction vector and $c \in \mathbb{R}$,
then $\text{boost}(c \cdot u) = c \cdot \text{boost}(u)$.
\end{theorem}
\begin{proof}
This follows from the linearity of the boost construction, proved above.
\end{proof}
\lean{Algebra.commuting\_family\_one\_parameter}

\begin{theorem}
If two boosts do not have parallel directions, they do not commute. Specifically, if there exist
$i, j$ with $v_i w_j \neq v_j w_i$, then $\text{boost}(v) \text{boost}(w) \neq \text{boost}(w) \text{boost}(v)$.
\end{theorem}
\begin{proof}
If they commuted, the previous theorem would force $v_i w_j = v_j w_i$ for all $i, j$, contradicting the hypothesis.
\end{proof}
\lean{Algebra.no\_two\_dimensional\_commuting\_family}

\begin{theorem}[Boost Drift]
Applying a boost to the distinguished basis vector returns the direction on the spatial block.
$\text{boost}(v)_{i+1, 0} = v_i$.
\end{theorem}
\begin{proof}
Immediate from the definition of the boost matrix.
\end{proof}
\lean{Algebra.boost\_drift}

\begin{theorem}[Drifts from Commuting Boosts are Proportional]
If $u$ and $v$ are commuting boost directions with $u_{i_0} \neq 0$, then there exists
a scalar $c$ such that $v_i = c \cdot u_i$ for all $i$.
\end{theorem}
\begin{proof}
By the maximal abelian subalgebra theorem, $v = \frac{v_{i_0}}{u_{i_0}} u$, so $c = \frac{v_{i_0}}{u_{i_0}}$ works.
\end{proof}
\lean{Algebra.commuting\_drifts\_proportional}

\begin{theorem}
The spatial block of a boost is zero. For any $i, j \in \{0, \ldots, k-1\}$,
$\text{boost}(v)_{i+1, j+1} = 0$.
\end{theorem}
\begin{proof}
This follows from the definition of the boost matrix.
\end{proof}
\lean{Algebra.boost\_spatial\_block\_zero}

\begin{theorem}
A boost is zero if and only if its direction is zero: $\text{boost}(v) = 0 \iff v = 0$.
\end{theorem}
\begin{proof}
If $\text{boost}(v) = 0$, then the $0, i+1$ entry is zero, i.e., $v_i = 0$ for all $i$, so $v = 0$.
Conversely, if $v = 0$, the matrix is identically zero.
\end{proof}
\lean{Algebra.boost\_eq\_zero\_iff}

\subsection{Dimension}

The distinction between the algebra and the space it acts on resolves an apparent degeneracy.
The algebra has dimension six while the space has four directions; these are not the same.
Once the algebra is named, the dimension of its defining representation is fixed, so $n$ becomes
dependent on the choice of algebra. The free parameter is now $k$ (spatial dimensions),
and a dimensional argument—requiring the generated rotations to carry labels—fixes $k = 3$ uniquely.

\begin{definition}
Twice the dimension of the rotation sector $\Lambda^2 p$ (with $k = m+1$ spatial directions):
$\text{rotDim2}(m) = (m+1) \cdot m$.
\end{definition}
\lean{Dimension.rotDim2}

\begin{definition}
Twice the dimension of the scaling sector $p$:
$\text{scaleDim2}(m) = 2(m+1)$.
\end{definition}
\lean{Dimension.scaleDim2}

\begin{definition}
Twice the dimension of the whole algebra $\mathrm{so}(k, 1)$ (with $k = m+1$):
$\text{algDim2}(m) = (m+1)(m+2)$.
\end{definition}
\lean{Dimension.algDim2}

\begin{definition}
Twice the dimension of the defining representation ($n = k + 1 = m + 2$):
$\text{repDim2}(m) = 2(m+2)$.
\end{definition}
\lean{Dimension.repDim2}

\begin{theorem}[Algebra Exceeds Representation]
The algebra is strictly larger than the space it acts on. Specifically,
$\text{repDim2}(2) < \text{algDim2}(2)$, i.e., $8 < 12$. The Lorentz algebra $\mathrm{so}(3,1)$
has dimension $6$, while spacetime has dimension $4$.
\end{theorem}
\begin{proof}
Direct calculation: $\text{repDim2}(2) = 2 \cdot 4 = 8$ and $\text{algDim2}(2) = 3 \cdot 4 = 12$.
\end{proof}
\lean{Dimension.algebra\_bigger\_than\_directions}

\begin{theorem}
Once the algebra is named, the representation dimension is determined. $\text{repDim2}(m) = 2((m+1) + 1)$ for all $m$.
\end{theorem}
\begin{proof}
Immediate from the definition.
\end{proof}
\lean{Dimension.rep\_dim\_determined}

\begin{theorem}[First Matching Condition]
The rotation sector has the same dimension as the scaling sector if and only if $m = 2$
(i.e., $k = 3$, $n = 4$). This is the condition $\text{rotDim2}(m) = \text{scaleDim2}(m)$,
which simplifies to $k(k-1)/2 = k$, solved uniquely by $k = 3$.
\end{theorem}
\begin{proof}
Setting $(m+1)m = 2(m+1)$ and dividing by $m+1$ gives $m = 2$.
\end{proof}
\lean{Dimension.rot\_matches\_scale\_iff}

\begin{theorem}[Second Matching Condition]
The algebra has the same dimension as its defining representation if and only if $m = 1$
(i.e., $k = 2$, $n = 3$). This is the condition $\text{algDim2}(m) = \text{repDim2}(m)$,
which simplifies to $k(k+1)/2 = k + 1$, solved uniquely by $k = 2$.
\end{theorem}
\begin{proof}
Setting $(m+1)(m+2) = 2(m+2)$ and dividing by $m+2$ gives $m+1 = 2$, so $m = 1$.
\end{proof}
\lean{Dimension.alg\_matches\_rep\_iff}

\begin{theorem}[The Two Conditions Disagree]
The two matching conditions give different answers:
\begin{itemize}
  \item At $m = 2$: $\text{rotDim2}(2) = \text{scaleDim2}(2) = 6$, but $\text{algDim2}(2) = 12 \neq 8 = \text{repDim2}(2)$.
  \item At $m = 1$: $\text{algDim2}(1) = \text{repDim2}(1) = 6$, but $\text{rotDim2}(1) = 2 \neq 4 = \text{scaleDim2}(1)$.
\end{itemize}
The statement ``require the structure to match itself'' is not well-posed until one specifies
which index set must match which. The framework chooses the first condition (rotation matches scale)
because A4 labels the scaling directions, and the coupling generates rotations that require those labels.
\end{theorem}
\begin{proof}
Direct calculation for both cases.
\end{proof}
\lean{Dimension.two\_matching\_conditions\_differ}

\begin{theorem}[Rescaling Invariance]
Rescaling the pattern rescales its rotational label by the square of the scaling factor.
For $c \in \mathbb{R}$ and $v, w \in \mathbb{R}^k$,
\[\text{wedge}(c \cdot v, c \cdot w)_{ij} = c^2 \cdot \text{wedge}(v, w)_{ij}.\]
\end{theorem}
\begin{proof}
By definition, $\text{wedge}(c \cdot v, c \cdot w)_{ij} = (cv_i)(cw_j) - (cv_j)(cw_i) = c^2(v_i w_j - v_j w_i)$.
\end{proof}
\lean{Dimension.wedge\_smul\_both}

\begin{theorem}[Ratio Invariance]
All ratios of rotational labels are invariant under rescaling the pattern. If $c \neq 0$ and $\text{wedge}(v,w)_{i'j'} \neq 0$, then
\[\frac{\text{wedge}(c \cdot v, c \cdot w)_{ij}}{\text{wedge}(c \cdot v, c \cdot w)_{i'j'}} = \frac{\text{wedge}(v, w)_{ij}}{\text{wedge}(v, w)_{i'j'}}.\]
Ratios are what the framework predicts: the conversion table, the gravitational numbers, and the dispersion.
\end{theorem}
\begin{proof}
By the previous theorem, $\text{wedge}(c \cdot v, c \cdot w)_{ij} = c^2 \text{wedge}(v, w)_{ij}$ and similarly for the denominator.
Dividing gives $\frac{c^2 \text{wedge}(v, w)_{ij}}{c^2 \text{wedge}(v, w)_{i'j'}} = \frac{\text{wedge}(v, w)_{ij}}{\text{wedge}(v, w)_{i'j'}}$.
\end{proof}
\lean{Dimension.wedge\_ratio\_invariant}

\begin{theorem}[Only Ratios are Fixed]
The framework fixes every prediction up to one overall factor, which it cannot determine because
there is nothing left to express it in. That factor is the unit. A dimensionless theory must have exactly one,
and this framework has exactly one—the minimum possible.
\end{theorem}
\begin{proof}
All physical quantities are ratios of wedges, which are invariant under simultaneous rescaling of the pattern.
No absolute scale of the pattern appears in any prediction; only their relative sizes matter.
\end{proof}
\lean{Dimension.only\_ratios\_are\_fixed}

\begin{definition}
The normalization $\alpha$ relating the threshold density $\rho$ to the rank-one parameter $k$:
\[\alpha(\rho, k) = \frac{k - 1}{\rho}.\]
This is a conversion relation—one equation, one unknown—not a prediction, and it requires the assumption
that defects are uniformly distributed in the symmetric space's volume.
\end{definition}
\lean{Dimension.alphaFromDensity}

\begin{theorem}[Density-Normalization Relation]
For nonzero density $\rho$, the product $\alpha(\rho, k) \cdot \rho$ equals $k - 1$:
\[\alpha(\rho, k) \cdot \rho = k - 1.\]
\end{theorem}
\begin{proof}
$\frac{k-1}{\rho} \cdot \rho = k - 1$ by definition.
\end{proof}
\lean{Dimension.density\_normalization\_relation}

\begin{theorem}
At $k = 3$ with the observed density $\rho$, the relation is $\alpha(\rho, 3) \cdot \rho = 2$.
\end{theorem}
\begin{proof}
Substituting $k = 3$ into the density-normalization relation gives $\alpha(\rho, 3) \cdot \rho = 3 - 1 = 2$.
\end{proof}
\lean{Dimension.normalization\_returns\_input}

\subsection{Slice}

Label proliferation with extension of scale range is unavoidable—every time a threshold is crossed,
new structures become resolvable and new labels are required. But this is not evidence of hidden structure;
it is the price of describing a wide range in a single formalism. The rotational quantity carried by the framework
is the wedge (exterior product) of two scale directions, and rotation exists only in the presence of
direction variation. The question "why three spatial dimensions" has a precise answer: $k = 3$ is the unique
dimension in which the generated rotations carry direction labels.

\begin{definition}
The content newly resolvable when extending the scale range from $t$ to $t'$ is the set of structures
whose thresholds lie in the gap:
\[\text{newContent}(\mu, t, t') = \{j : \mu_j \leq t' \text{ and } \mu_j > t\}.\]
\end{definition}
\lean{Slice.newContent}

\begin{theorem}
Membership in $\text{newContent}(\mu, t, t')$ is equivalent to: $\mu_j \leq t'$ and $\mu_j \not\leq t$.
\end{theorem}
\begin{proof}
Immediate from the definition.
\end{proof}
\lean{Slice.mem\_newContent}

\begin{theorem}[Extension Forces New Labels]
If any structure has its threshold strictly between $t$ and $t'$, the extended description must carry
a label the local one does not need. Specifically, if $t < \mu_j \leq t'$ for some $j$, then
$\text{newContent}(\mu, t, t')$ is nonempty.
\end{theorem}
\begin{proof}
The structure $j$ satisfies $\mu_j \leq t'$ (resolvable at the top) and $\mu_j > t$ (not resolvable at the bottom),
so $j \in \text{newContent}(\mu, t, t')$.
\end{proof}
\lean{Slice.newContent\_nonempty\_of\_threshold}

\begin{theorem}[Excess is Monotone in Range]
Widening the scale window can only add new labels, never remove them. If $t_1 \leq t_2$ and $t' \leq t''$, then
\[\text{newContent}(\mu, t_2, t') \subseteq \text{newContent}(\mu, t_1, t'').\]
\end{theorem}
\begin{proof}
If $j \in \text{newContent}(\mu, t_2, t')$, then $\mu_j \leq t' \leq t''$ and $\mu_j > t_2 \geq t_1$, so $j \in \text{newContent}(\mu, t_1, t'')$.
\end{proof}
\lean{Slice.newContent\_mono}

\begin{theorem}[Threshold-Free Ranges Cost Nothing]
If there is no structure with threshold in the range $(t, t']$, then the extended description is exact:
no new labels are required.
\end{theorem}
\begin{proof}
If $\text{newContent}(\mu, t, t')$ is empty, then no structure $j$ satisfies both $\mu_j \leq t'$ and $\mu_j > t$,
which means every structure either is not resolvable at $t'$ or is already resolvable at $t$.
\end{proof}
\lean{Slice.newContent\_empty\_of\_no\_threshold}

\begin{theorem}[No Universal Labelling]
No fixed label set works at every scale. Given any threshold function $\mu$ with $\mu_j \to \infty$,
for any scale $t$ there exists a higher scale $t'$ such that new content appears.
\end{theorem}
\begin{proof}
By the unboundedness assumption, there exists a structure $j$ with $\mu_j > t$. Setting $t' = \mu_j$ gives
$t < t' \leq t'$ and $\mu_j > t$, so $j \in \text{newContent}(\mu, t, t')$.
\end{proof}
\lean{Slice.no\_finite\_universal\_labelling}

\begin{theorem}[Label Count Measures Range]
The number of labels a formalism carries measures the log-range it covers, not the depth of the physics.
A description at a single scale needs only the content of that scale; a description across a range
must carry labels for all content across that range.
\end{theorem}
\begin{proof}
This is the collective statement of the previous theorems: new labels accumulate monotonically with range,
and any fixed labelling eventually becomes insufficient.
\end{proof}
\lean{Slice.label\_count\_measures\_range}

\begin{definition}
The \textbf{wedge} (exterior product) of two scale directions $v, w \in \mathbb{R}^k$ is the antisymmetric $k \times k$ matrix
\[\text{wedge}(v, w)_{ij} = v_i w_j - v_j w_i.\]
This is the rotational label generated by two scale directions.
\end{definition}
\lean{Slice.wedge}

\begin{theorem}[Wedge Equals Bracket]
The wedge is exactly what the coupling generates. For boost matrices,
\[\text{wedge}(v, w)_{ij} = [\text{boost}(v), \text{boost}(w)]_{i+1,j+1},\]
where the right side is the $(i+1, j+1)$ entry of the commutator.
\end{theorem}
\begin{proof}
This follows from Algebra.boost\_bracket\_spatial, which computed the bracket entry by entry.
\end{proof}
\lean{Slice.wedge\_eq\_bracket}

\begin{theorem}[Wedge is Antisymmetric]
For all $i, j$,
\[\text{wedge}(v, w)_{ij} = -\text{wedge}(v, w)_{ji}.\]
\end{theorem}
\begin{proof}
$\text{wedge}(v, w)_{ij} = v_i w_j - v_j w_i = -(v_j w_i - v_i w_j) = -\text{wedge}(v, w)_{ji}$.
\end{proof}
\lean{Slice.wedge\_antisymm}

\begin{theorem}[Wedge is Bilinear]
The wedge is bilinear: $\text{wedge}(u + v, w) = \text{wedge}(u, w) + \text{wedge}(v, w)$.
\end{theorem}
\begin{proof}
$\text{wedge}(u + v, w)_{ij} = (u_i + v_i) w_j - (u_j + v_j) w_i = (u_i w_j - u_j w_i) + (v_i w_j - v_j w_i)$.
\end{proof}
\lean{Slice.wedge\_bilinear\_left}

\begin{theorem}[Wedge Vanishes Iff Parallel]
Two directions generate no rotation if and only if they are parallel (componentwise proportional).
$\text{wedge}(v, w) = 0$ if and only if $v_i w_j = v_j w_i$ for all $i, j$.
\end{theorem}
\begin{proof}
Each entry $\text{wedge}(v, w)_{ij} = v_i w_j - v_j w_i$ vanishes iff $v_i w_j = v_j w_i$.
\end{proof}
\lean{Slice.wedge\_eq\_zero\_iff\_parallel}

\begin{theorem}[Single Direction Carries No Rotation]
A homogeneous pattern (one direction, possibly with varying magnitude) carries no rotational label.
$\text{wedge}(v, v)_{ij} = 0$ for all $i, j$.
\end{theorem}
\begin{proof}
$\text{wedge}(v, v)_{ij} = v_i v_j - v_j v_i = 0$.
\end{proof}
\lean{Slice.homogeneous\_carries\_no\_rotation}

\begin{theorem}[Nonzero Rotation Requires Two Non-Parallel Directions]
If a wedge is nonzero, no scalar $c$ satisfies $w = c \cdot v$.
\end{theorem}
\begin{proof}
If $w = c \cdot v$, then $\text{wedge}(v, w)_{ij} = v_i (cv_j) - v_j (cv_i) = 0$ for all $i, j$.
Contrapositive: if $\text{wedge}(v, w)_{ij} \neq 0$ for some $i, j$, then $w \neq c \cdot v$ for any $c$.
\end{proof}
\lean{Slice.rotation\_requires\_biaxial}

\begin{definition}
The number of independent wedges of $k$ directions (equivalently, the dimension of $\Lambda^2 \mathbb{R}^k$):
\[\text{wedgeDim}(m) = (m+1) \cdot m \quad \text{(with $k = m+1$)}.\]
\end{definition}
\lean{Slice.wedgeDim}

\begin{definition}
The number of directions, written to match the scaling of $\text{wedgeDim}$:
\[\text{dirDim}(m) = 2(m+1).\]
\end{definition}
\lean{Slice.dirDim}

\begin{theorem}[Unique Matching Dimension]
The wedges are as numerous as the directions if and only if $m = 2$, i.e., $k = 3$.
The equation $k(k-1)/2 = k$ has unique positive solution $k = 3$.
\end{theorem}
\begin{proof}
Setting $\text{wedgeDim}(m) = \text{dirDim}(m)$ gives $(m+1)m = 2(m+1)$. Dividing by $m+1$ (which is positive) yields $m = 2$.
\end{proof}
\lean{Slice.wedge\_dim\_eq\_iff}

\begin{theorem}[Why Three Dimensions]
Three spatial dimensions is the unique dimension in which the coupling closes on the labels the axioms provide.
Above $k = 3$ there are more wedges than directions, so some generated rotation is unlabelled; below $k = 3$
the rotation sector degenerates. This is the unique $m$ satisfying $\text{wedgeDim}(m) = \text{dirDim}(m)$.

The arithmetic is a theorem. The physics—that the framework \emph{requires} its generated rotations to be
direction-labelled rather than merely permitting some to go unlabelled—is a judgement registered as such,
on the same footing as the rank-one identification.
\end{theorem}
\begin{proof}
The existence of $m = 2$ is direct calculation. Uniqueness follows from the fact that $m = 2$ is the unique
solution to $(m+1)m = 2(m+1)$.
\end{proof}
\lean{Slice.three\_is\_unique}

\begin{theorem}[Higher Dimensions Have More Wedges]
For $m > 2$ (i.e., $k > 3$), the wedges outnumber the directions: $\text{dirDim}(m) < \text{wedgeDim}(m)$.
\end{theorem}
\begin{proof}
For $m > 2$, we have $(m+1)m > 2(m+1)$ because $m > 2$.
\end{proof}
\lean{Slice.more\_wedges\_than\_directions}

\begin{theorem}[Lower Dimensions Have Fewer Wedges]
For $m < 2$ (i.e., $k < 3$), the wedges are fewer than the directions: $\text{wedgeDim}(m) < \text{dirDim}(m)$.
\end{theorem}
\begin{proof}
By direct case analysis: for $m = 0$, $\text{wedgeDim}(0) = 0 < 2 = \text{dirDim}(0)$; for $m = 1$, $\text{wedgeDim}(1) = 2 < 4 = \text{dirDim}(1)$.
\end{proof}
\lean{Slice.fewer\_wedges\_than\_directions}

\subsection{Axes}

The resolution of the apparent dilemma between rank-one and angular momentum is a misreading.
The theorem $\text{wedge}(v, v) = 0$ says a homogeneous (constant) pattern carries no rotation, not that a uniaxial pattern
carries none. The pattern is a single vector, and rotation arises only from its *variation* in space.
A topologically charged defect—one whose pattern winds—always carries angular momentum.
Real rank one forces the pattern to be a single vector, so biaxial or higher-order patterns are not available.

\begin{theorem}[Rotation Requires Variation]
Nonzero rotational label requires two non-parallel values of the pattern.
\end{theorem}
\begin{proof}
At any one place, the pattern is a single vector, and its wedge with itself vanishes. Rotation can arise
only between the pattern's values at different points.
\end{proof}
\lean{Axes.rotation\_needs\_variation}

\begin{theorem}[Parallel Patterns Carry No Rotation]
A pattern whose values are all parallel (a multiple of one fixed direction at each point) generates
no rotation anywhere, because every wedge vanishes.
\end{theorem}
\begin{proof}
If $v_1, v_2, \ldots$ are all multiples of one direction $u$, then $\text{wedge}(f_m \cdot u, f_l \cdot u) = 0$ for any scalars $f_m, f_l$.
\end{proof}
\lean{Axes.combable\_no\_rotation}

\begin{theorem}[Wound Defect Carries Rotation]
A pattern taking two non-parallel values generates nonzero rotational label. If $v_i w_j \neq v_j w_i$
for some $i, j$, then $\text{wedge}(v, w)_{ij} \neq 0$.
\end{theorem}
\begin{proof}
By definition, $\text{wedge}(v, w)_{ij} = v_i w_j - v_j w_i$, which is nonzero by hypothesis.
\end{proof}
\lean{Axes.wound\_carries\_rotation}

\begin{theorem}[Rotationless Configuration is Trivial]
If the rotational label vanishes everywhere, the configuration is topologically trivial.
The framework forbids a charged defect with zero angular momentum.
\end{theorem}
\begin{proof}
If $\text{wedge}(v, w)_{ij} = 0$ for all $i, j$, then all values of the pattern are parallel, so the configuration
is combable and deforms to a constant.
\end{proof}
\lean{Axes.rotationless\_is\_trivial}

\begin{theorem}[Rotation Iff Variation]
Rotation exists if and only if the pattern's values are not everywhere parallel.
\end{theorem}
\begin{proof}
The equivalence is immediate from the definitions of $\text{wedge}$ and the notion of parallel vectors.
\end{proof}
\lean{Axes.rotation\_iff\_variation}

\begin{theorem}[No Second Axis]
Real rank one forces the pattern to be a single vector. Any other scaling is either a multiple of it
(the same axis) or has nonzero wedge with it (showing up as rotation, not a second axis).
If $v$ has a nonzero component and $v_i w_j = v_j w_i$ for all $i, j$, then $w$ is a multiple of $v$:
$w = \frac{w_{i_0}}{v_{i_0}} v$.
\end{theorem}
\begin{proof}
From $v_{i_0} w_i = v_i w_{i_0}$ for all $i$, we get $w_i = \frac{w_{i_0}}{v_{i_0}} v_i$.
\end{proof}
\lean{Axes.no\_second\_axis}

\begin{theorem}[Axis or Variation]
Given the pattern $v$ with a nonzero component, every other scaling $w$ is either parallel to $v$ or has
nonzero wedge with it.
\end{theorem}
\begin{proof}
By the wedge criterion, either all wedges vanish (parallel case) or some wedge is nonzero (variation case).
\end{proof}
\lean{Axes.axis\_or\_variation}

\begin{theorem}[Stabilizer is Positive-Dimensional]
The rotation generator in the plane transverse to the pattern annihilates it. This is the signature of
a uniaxial order parameter: a continuous stabilizer. A biaxial or fully anisotropic pattern would have
a discrete stabilizer, which does not support point defects.

For the concrete $3 \times 3$ rotation $J$ (in the $1$--$2$ plane) applied to the distinguished direction
$(1, 0, 0)$ in the standard basis, the result is zero: $J \cdot (1, 0, 0) = 0$. And $J$ is nonzero.
\end{theorem}
\begin{proof}
Computing $J \cdot (1, 0, 0) = (0, 0, 0) - (0, 0, 0) = 0$ while $J \neq 0$.
\end{proof}
\lean{Axes.axis\_stabilizer\_nontrivial}

\begin{theorem}[Multiplet Flatness Requires Per-Direction Scalings]
The theorem that ``multiplets are internally flat'' required a scaling operator per direction.
Real rank one provides only one vector and one scaling, so that physical reading is withdrawn.
What stands is that a homogeneous pattern generates no rotation, which is a statement about constancy,
not about the relation between directions.
\end{theorem}
\begin{proof}
This is immediate from the fact that $\text{wedge}(v, v) = 0$ for any vector $v$.
\end{proof}
\lean{Axes.multiplet\_flatness\_needs\_per\_direction\_scalings}

